% ===== 原文 pp.87--92 =====
% [原文 p.87]
\subsubsection{离散挠率（discrete torsion）}

上述规范化程序有一个轻微的推广：在 (4.139) 中出现的求和或积分上添加系数
\begin{equation}
\begin{aligned}
Z_{\mathcal{T}/G^{(p)}}&\propto\sum_{[B_{p+1}]}\mathcal{I}_d[B_{p+1}]\,Z_{\mathcal{T}}[B_{p+1}],\qquad G^{(p)}\ \text{离散},\\[2pt]
Z_{\mathcal{T}/G^{(p)}}&\propto\int\mathcal{D}[B_{p+1}]\,\mathcal{I}_d[B_{p+1}]\,Z_{\mathcal{T}}[B_{p+1}],\qquad G^{(p)}\ \text{连续}.
\end{aligned}
\tag{4.158}
\end{equation}
这些系数本身必须是规范不变的
\begin{equation}
\mathcal{I}_d[B^g_{p+1}]=\mathcal{I}_d[B_{p+1}]\,.
\tag{4.159}
\end{equation}
原则上，这些系数 $\mathcal{I}_d[B_{p+1}]$ 可以是任何具有非反常 $G^{(p)}$ 对称性的 $d$ 维理论 $\mathcal{T}'$ 的配分函数。在这种情况下，上述规范化程序相当于把 $\mathcal{T}$ 与一份解耦的 $\mathcal{T}'$ 副本堆叠（stacking），得到总对称性 $G^{(p)}\times G^{(p)}$，然后再规范化对角 $G^{(p)}$ 对称性。

然而，$\mathcal{I}_d$ 是 $d$ 维 $G^{(p)}$ SPT 相的情形特别有意思，因为它只是提供了另一种规范化 $\mathcal{T}$ 的 $G^{(p)}$ $p$-形式对称性的方式。这是因为 SPT 相（在忘掉对称性后）是一个平凡理论，因此与 SPT 相堆叠不会给 $\mathcal{T}$ 增添任何新的自由度。在这种情形下，当 $\mathcal{I}_d$ 是 SPT 相时，规范化程序 (4.158) 被称为带由 SPT 相 $\mathcal{I}_d$ 给出的\textbf{离散挠率}（discrete torsion）的规范化。

% [Tutor Note: p.87 式 (4.157)、(4.158)——(4.158) 只是在 (4.139) 的每一项前乘一个规范不变的系数。若 $\mathcal{I}_d$ 取任意带非反常 $G^{(p)}$ 对称性的理论 $\mathcal{T}'$ 的配分函数，则等价于先把 $\mathcal{T}$ 与 $\mathcal{T}'$ 堆叠（得到 $G^{(p)}\times G^{(p)}$ 对称性）再规范化对角子群；核对：$G^{(p)}$ 离散时两个因子的求和指标都是同一个 $[B_{p+1}]$，故 (4.158) 的写法与「对角规范化」自洽。只有取 $\mathcal{I}_d$ 为 SPT 相时堆叠才不引入新自由度，这才是原文说「特别有意思」的原因。]

\begin{examplebox}{例 4.11：从 $\SU(2)$ 到 $\SO(3)_-$}
考虑 4d 纯 $\SU(2)$ Yang--Mills 理论，它具有 1-形式对称性
\begin{equation}
G^{(1)}=\Z_2\,.
\tag{4.160}
\end{equation}
4d 中存在 $\Z_2$ 1-形式对称性的一个 SPT 相，它在背景场存在时的配分函数为
\begin{equation}
\mathcal{I}_4[B_2]=\exp\left(\pi\ii\int\frac{\mathcal{P}(B_2)}{2}\right).
\tag{4.161}
\end{equation}
用离散挠率 $\mathcal{I}_4$ 对 $\SU(2)$ 理论的 $\Z_2$ 1-形式对称性作规范化，得到的规范理论具有 $\SO(3)$ 规范群，并在作用量中带有项 (3.193)。这恰好就是 3.4.1 节中讨论的 4d 纯 $\SO(3)_-$ Yang--Mills 理论，它携带一个离散 theta 角。于是我们得知
\begin{equation}
\left(\mathcal{T}_{\SU(2)}\times\mathcal{I}_4\right)/\Z_2^{(1)}=\mathcal{T}_{\SO(3)_-}.
\tag{4.162}
\end{equation}
\end{examplebox}

% [Tutor Note: p.87 式 (4.161)——$\mathcal{P}$ 是 Pontryagin 平方（见式 (3.193) 附近的讨论）。可核对之处：对 $\Z_2$-值 2-上链 $B_2$，$\mathcal{P}(B_2)$ 取偶值，因此 $\mathcal{P}(B_2)/2$ 在模 1 意义下良定义，$\exp(\pi\ii\int \mathcal{P}(B_2)/2)$ 才在背景规范变换下不变；若漏掉分母 2，相位会变成 $\pi\ii\int\mathcal{P}$，对 $\Z_2$ 值仍不变但会与 (3.193) 的归一化差一档，导致 (4.162) 中的 $\SO(3)_-$ 离散 theta 角不对。]

% [Tutor Note: p.87 式 (4.162)——$\SO(3)_-$ 的下标「$-$」标记的是 4d 纯 $\SO(3)$ Yang--Mills 中作用量含离散 theta 角项 (3.193) 的那一支（见 3.4.1 节；通常的 $\SO(3)$ 理论相应记作 $\SO(3)_+$）。从 $\SU(2)$ 出发加上离散挠率 $\mathcal{I}_4$ 恰好选出 $\SO(3)_-$ 这一支，这就是「离散挠率 $\leftrightarrow$ 离散 theta 角」的对应，也是 4.3.4 节讨论的先声。]

\subsubsection{规范化后得到的对偶对称性}

考虑规范化 (4.137)，它在规范化理论 $\mathcal{T}/G^{(p)}$ 中给出一个 $G^{(p)}$-值规范场 $b_{p+1}$。与任何规范场一样，我们可以考虑 $b_{p+1}$ 的 Wilson 算符 $\mathcal{W}_R$，它们由 $G^{(p)}$ 的不可约表示 $R$ 参数化。这些算符是 $\mathcal{T}/G^{(p)}$ 的 $(p+1)$ 维算符。

% [原文 p.88]
若 $G^{(p)}$ 是连续的，则 $\mathcal{W}_R$ 是非拓扑算符。然而，若 $G^{(p)}$ 是离散的，则由于此时 $b_{p+1}$ 必然是平坦的，$\mathcal{W}_R$ 是拓扑算符。既然 $\mathcal{W}_R$ 是拓扑的，它们必定在规范化理论中生成一个新的对称性。于是我们得到如下陈述：

\begin{statementbox}{陈述 4.5：规范化理论的对偶对称性}
对离散 \emph{阿贝尔} $G^{(p)}$ $p$-形式对称群规范化后得到的规范化理论 $\mathcal{T}/G^{(p)}$ 具有一个 $(d-p-2)$-形式对称性，它由 Pontryagin 对偶群
\begin{equation}
G^{(d-p-2)}=\widehat{G}^{(p)}\,.
\tag{4.163}
\end{equation}
给出。这类对称性称为\textbf{对偶对称性}（dual symmetries）或\textbf{量子对称性}（quantum symmetries），它们来自 $G^{(p)}$ 的规范化。
\end{statementbox}

注意上述陈述中要求 $G^{(p)}$ 是 阿贝尔的。对 $p>0$，$G^{(p)}$ 必须是 阿贝尔的；但对 $p=0$，$G^{(0)}$ 可以是非阿贝尔的。这带来如下有意思的推论：

\begin{notebox}{不可逆的表示对称性}
首先注意，对离散非阿贝尔群 $G^{(0)}$，总存在维数大于 1 的不可约表示 $R$。因此不可能找到 $G^{(0)}$ 的另一个表示 $R^*$ 使得
\begin{equation}
R\otimes R^*=\id\,,
\tag{4.164}
\end{equation}
其中右边是平凡表示。这是因为表示的维数在张量积下相乘。

于是，在规范化理论 $\mathcal{T}/G^{(0)}$ 中，我们有一些\textbf{非可逆}的拓扑 Wilson 线算符。由非可逆拓扑算符生成的对称性称为不可逆对称性，近来已成为研究的主要焦点 [56--59,122--177]。在上述情形下，人们说 $\mathcal{T}/G^{(0)}$ 具有一个
\begin{equation}
\mathrm{Rep}(G^{(0)})\,,
\tag{4.165}
\end{equation}
非可逆 $(d-2)$-形式对称性。
\end{notebox}

% [Tutor Note: p.88 式 (4.163)——Pontryagin 对偶群 $\widehat{G}^{(p)}=\mathrm{Hom}(G^{(p)},\U(1))$。原文强调 $G^{(p)}$ 必须 阿贝尔：$p>0$ 时 $(p+1)$-形式算符的融合可交换，高阶形式对称群自动 阿贝尔，故该限制无损；$p=0$ 时非阿贝尔群仍可商，但对偶不再是群的 Pontryagin 对偶，而要用表示范畴 $\mathrm{Rep}(G^{(0)})$，这正是 (4.165) 出现的原因。核对：$\widehat{\Z_n}=\Z_n$，故离散 $\Z_n$ 1-形式对称性规范化后得到对偶 $\Z_n$ $(d-3)$-形式对称性。]

% [Tutor Note: p.88 式 (4.164)——不可逆性的判据是「不存在逆算符」，等价于 $R\otimes R^*$ 不等于平凡表示。维数论证的可用锚点：取 $G^{(0)}=S_3$ 的 2 维不可约表示 $R$，则 $R\otimes R^*$ 的维数为 $2\times2=4$，不可能是 1 维平凡表示，故相应的 Wilson 线不可逆。英文原词：non-invertible symmetry（不可逆对称性）。]

% [原文 p.89]
\para{对偶的对偶。}现在考虑规范化对偶对称性 (4.163)。结果表明，所得的规范化理论就是原来的理论 $\mathcal{T}$，即
\begin{equation}
\left(\mathcal{T}/G^{(p)}\right)/G^{(d-p-2)}=\mathcal{T}\,,
\tag{4.166}
\end{equation}
而对偶对称性 $G^{(d-p-2)}$ 的对偶就是原来的对称性 $G^{(p)}$。

这可以通过注意到下述事实很快看出：开启对偶对称性的背景场 $B_{d-p-1}$ 会添加耦合
\begin{equation}
2\pi\int B_{d-p-1}\cup_\eta b_{p+1}\,.
\tag{4.167}
\end{equation}
对 $B_{d-p-1}$ 作规范化，并加入与新对偶 $p$-形式对称性背景场 $B^{\text{new}}_{p+1}$ 的耦合后，我们有
\begin{equation}
2\pi\int b_{d-p-1}\cup_\eta b_{p+1}+B^{\text{new}}_{p+1}\cup_\eta b_{d-p-1}\,.
\tag{4.168}
\end{equation}
积掉 $b_{d-p-1}$ 会令
\begin{equation}
b_{p+1}=(-1)^p B^{\text{new}}_{p+1}\,.
\tag{4.169}
\end{equation}
这样我们确实（差一个符号地）回到了原来的理论 $\mathcal{T}$，从而 (4.166) 得到说明，并且发现新的对偶对称性就是原来的 $G^{(p)}$ 对称性。

% [Tutor Note: p.89 式 (4.167)--(4.169)——$(-1)^p$ 的来源是 Koszul 符号：$b_{p+1}$ 与 $b_{d-p-1}$ 交换时要产生 $(-1)^{(p+1)(d-p-1)}$，而在 $d-p-1\equiv 1\pmod 2$ 的奇次数情形它约化为 $(-1)^p$；这也是 (4.94) 中 $s(p)=(-1)^p$ 的同源符号。核对：$p=1,d=4$ 时 (4.169) 给出 $b_2=-B^{\text{new}}_2$，与原文「up to a sign」的措辞一致。]

\para{非阿贝尔规范理论中电 1-形式对称性的规范化。}正如 4.3.1 节所见，规范化规范理论的电 1-形式对称性会缩小规范群。更准确地说，它对规范群实现一个离散商。较小的规范群具有较大的磁 $(d-3)$-形式对称性。这些新的磁对称性可以理解为由 1-形式规范化产生的对偶对称性。在叙述一般结果之前，我们先看一个例子。

\begin{examplebox}{例 4.12：从 $\SU(2)$ 到 $\SO(3)$}
回忆纯 $\SU(2)$ Yang--Mills 理论 $\Z_2$ 电 1-形式对称性的背景场被等同于障碍类
\begin{equation}
B_2=w_2\,,
\tag{4.170}
\end{equation}
它对应把 $\SO(3)$ 丛提升为 $\SU(2)$ 丛的障碍。对 $\Z_2$ 1-形式对称性作规范化后，我们得到纯 $\SO(3)$ Yang--Mills 理论。对偶的 $\Z_2$ $(d-3)$-形式对称性由拓扑算符
\begin{equation}
\exp\left(\ii\pi\int b_2\right)=\exp\left(\ii\pi\int w_2\right).
\tag{4.171}
\end{equation}
生成。
% [原文 p.90]
现在回忆 (3.113) 可知，上式右端正是生成 $\SO(3)$ 理论磁 $(d-3)$-形式对称性的拓扑算符。因此，对偶 $(d-3)$-形式对称性正是我们一直在研究的磁 $(d-3)$-形式对称性。
\end{examplebox}

% [Tutor Note: p.89 式 (4.170)--(4.171)——核对符号：$B_2$ 是背景场，$w_2$ 是丛的障碍类，二者在纯 $\SU(2)$ 理论中等同，故 (4.171) 把 $\int b_2$ 与 $\int w_2$ 写成一个算子；规范化后 $w_2$ 变为动力学的 $b_2$。与 (4.142) $w_2=B_2^e$ 对照可确认上标 $e$ 在例 4.12 中被省略。]

更一般地，对非阿贝尔规范理论的电 1-形式对称性作规范化，我们有结果 (4.147)。$\mathcal{T}_{\mathcal{G}}$ 与 $\mathcal{T}_{\mathcal{G}/H^{(1)}}$ 的磁 $(d-3)$-形式对称性分别为
\begin{equation}
G^{(d-3)}(\mathcal{T}_{\mathcal{G}})=\widehat{\mathcal{Z}},\qquad G^{(d-3)}(\mathcal{T}_{\mathcal{G}/H^{(1)}})=\widehat{\mathcal{Y}}\,,
\tag{4.172}
\end{equation}
其中 $\mathcal{Z}$ 与 $\mathcal{Y}$ 由
\begin{equation}
\mathcal{G}=G/\mathcal{Z},\qquad \mathcal{G}/H^{(1)}=G/\mathcal{Y}\,,
\tag{4.173}
\end{equation}
定义（这里用的是单连通群 $G$）。这两个群 $\mathcal{Z}$ 与 $\mathcal{Y}$ 由短正合列
\begin{equation}
0\to\mathcal{Z}\to\mathcal{Y}\to H^{(1)}\to 0
\tag{4.174}
\end{equation}
相联系，而其 Pontryagin 对偶短正合列把磁对称性 (4.172) 联系为
\begin{equation}
0\to\widehat{H}^{(1)}\to\widehat{\mathcal{Y}}\to\widehat{\mathcal{Z}}\to 0\,.
\tag{4.175}
\end{equation}
$G^{(d-3)}(\mathcal{T}_{\mathcal{G}/H^{(1)}})=\widehat{\mathcal{Y}}$ 的子群 $\widehat{H}^{(1)}$ 就是由规范化 $H^{(1)}$ 1-形式对称性产生的对偶 $(d-3)$-形式对称性。

% [Tutor Note: p.90 式 (4.174)--(4.175)——取 Pontryagin 对偶把短正合列反向：$0\to\mathcal{Z}\to\mathcal{Y}\to H^{(1)}\to 0$ 的对偶是 $0\to\widehat{H}^{(1)}\to\widehat{\mathcal{Y}}\to\widehat{\mathcal{Z}}\to 0$（箭头全部反向）。核对锚点：对有限 阿贝尔群，$\widehat{\cdot}$ 是正合反变函子，故正合性保持；用 $\mathcal{Z}=\Z_2,\mathcal{Y}=\Z_4,H^{(1)}=\Z_2$ 检验，(4.175) 变成 $0\to\Z_2\to\Z_4\to\Z_2\to 0$，与 (4.126) 同形，自洽。]

\subsubsection{规范化交换扭曲扇区与带电扇区}

\para{扭曲变为带电。}在 $\mathcal{T}/G^{(p)}$ 的对偶对称性 $G^{(d-p-2)}$ 下带电的缺陷，处在 $\mathcal{T}$ 的对称性 $G^{(p)}$ 的扭曲扇区中。

\begin{definitionbox}{定义 4.7：扭曲扇区}
称一个余维 $(p+2)$ 算符 $\mathcal{O}$ 处于某个 $p$-形式对称性的扭曲扇区，若 $\mathcal{O}$ 是出现在与该 $p$-形式对称性相联的余维 $(p+1)$ 拓扑算符端点处的非真算符。
\end{definitionbox}

% [Tutor Note: p.90 定义 4.7——「扭曲扇区」（twisted sector）的判据有两个要素：算符是非真的，且它只作为某个拓扑算符的端点出现。注意与一般非真算符区分：本定义中的非真专门来自 $G^{(p)}$ 的拓扑算符端点（英文原词 twisted sector）。核对锚点：余维数必须是 $p+2$，因为相联的拓扑算符是余维 $p+1$，其端点比它多一维；若写成余维 $p+1$ 就把算符与其端点混为一谈。]

为看清这一点，考虑一个扭曲扇区算符 $\mathcal{O}$，它沿时空的一个子流形 $M_{d-p-2}$ 插入，位于拓扑算符
\begin{equation}
U_g\,,\qquad g\in G^{(p)}\,.
\tag{4.176}
\end{equation}
的端点。对 $G^{(p)}$ 作规范化后，拓扑算符 $U_g$ 实现的是规范变换，因而不作用于任何物理可观测量。因此，在 $\mathcal{T}/G^{(p)}$ 中算符 $U_g$ 可以被等同于恒等算符。人们常说 $U_g$ 在规范化后变得\textbf{不可见}（invisible）。于是 $\mathcal{O}$ 是规范化理论 $\mathcal{T}/G^{(p)}$ 的真余维 $(p+2)$ 算符。见图 18。

% [原文 p.91]
\begin{figure}[htbp]
\centering
\includegraphics[width=0.9\textwidth]{fig18.png}
\figtitle{18}{该图展示：处于 $G^{(p)}$ $p$-形式对称性扭曲扇区中的非真余维 $(p+2)$ 算符 $\mathcal{O}$，在对 $G^{(p)}$ 作规范化后变成真算符。严格地说，只有当 $\mathcal{O}$ 在 $G^{(p)}$ 下不带荷时这一点才成立，这将在后文讨论。}
\end{figure}

现在我们断言 $\mathcal{O}$ 在 $\mathcal{T}/G^{(p)}$ 的对偶 $G^{(d-p-2)}$ 对称性下带有电荷
\begin{equation}
g\in G^{(p)}=\widehat{\widehat{G}}^{(p)}=\widehat{G}^{(d-p-2)}\,,
\tag{4.177}
\end{equation}
这可以从下述事实看出：在 $\mathcal{T}$ 中，$\mathcal{O}(M_{d-p-2})$ 的插入会对背景场施加
\begin{equation}
\delta B_{p+1}=\delta_g(M_{d-p-2})\,,
\tag{4.178}
\end{equation}
这只是「$\mathcal{O}$ 处于 $g$-扭曲扇区」这一事实的另一种表述。

对 $G^{(p)}$ 作规范化后，(4.178) 变为
\begin{equation}
\delta b_{p+1}=\delta_g(M_{d-p-2})\,.
\tag{4.179}
\end{equation}
回忆 (2.105) 附近的讨论可知，这恰好意味着 $\mathcal{O}$ 在 $G^{(d-p-2)}$ 下带电荷 $g$。

% [Tutor Note: p.91 式 (4.178)--(4.179)——$\delta_g(M_{d-p-2})$ 指支撑在 $M_{d-p-2}$ 上、取值为 $g$ 的 delta 函数型 $(p+2)$-上链（即「插入一个带 $g$ 荷的算符」在背景场层面上的体现）；规范化前后只是把背景场 $B_{p+1}$ 换成动力学场 $b_{p+1}$，等式形式不变。核对方法：对 $\Z_2$ 1-形式对称性，$\delta_g$ 是 $M_{d-3}$ 上的 $\Z_2$ 值 3-上链，$\delta b_2=\delta_g$ 表示 $b_2$ 沿一个 2-球面围绕 $M_{d-3}$ 的积分为 $g$，与 (2.105) 的荷定义一致。]

\begin{examplebox}{例 4.13：从 $\SU(2)$ 到 $\SO(3)$：'t Hooft 线}
在 $d$ 维纯 $\SO(3)$ Yang--Mills 理论的 $\Z_2$ 磁 $(d-3)$-形式对称性下带电的算符，是那些诱导出不能提升为 $\SU(2)$ 磁单极子位形的 $\SO(3)$ 磁单极子位形的 't Hooft 算符。在纯 $\SU(2)$ Yang--Mills 理论中，这类 't Hooft 算符的 Dirac 弦是可见的，并且可以被等同于生成 $\SU(2)$ 理论 $\Z_2$ 1-形式对称性的拓扑余维-2 算符。因此，这类 't Hooft 算符处在 $\SU(2)$ 理论中电 1-形式对称性的扭曲扇区中。
\end{examplebox}

% [Tutor Note: p.91 例 4.13——这是图 18 的显式实例：同一族 't Hooft 线在 $\SU(2)$ 理论里「Dirac 弦可见」故是非真算符（扭曲扇区），在 $\SO(3)$ 理论里 Dirac 弦被规范化掉、从而变成真算符并带磁 $\Z_2$ 荷。可核对之处：Dirac 弦正是 (3.113) 的拓扑余维-2 算符，即 $\Z_2$ 1-形式对称性的生成元，与定义 4.7 中的「拓扑算符端点」吻合。]

\para{带电变为扭曲。}另一方面，在 $\mathcal{T}$ 的 $G^{(p)}$ 对称性下带电的算符，被等同于 $\mathcal{T}/G^{(p)}$ 的 $G^{(d-p-2)}$ 对称性的扭曲扇区算符。为看清这一点，我们从 $\mathcal{T}/G^{(p)}$ 中的扭曲扇区算符出发并对 $G^{(d-p-2)}$ 作规范化。正如 (4.166) 附近所讨论的，这会把我们带回理论 $\mathcal{T}$。按上述讨论，$\mathcal{T}/G^{(p)}$ 的扭曲扇区算符在这一过程中变成 $\mathcal{T}$ 的带电算符。

% [原文 p.92]
\begin{figure}[htbp]
\centering
\includegraphics[width=0.85\textwidth]{fig19.png}
\figtitle{19}{一个既处于 $G^{(p)}$ $p$-形式对称性的扭曲扇区、又在该对称性下带电的算符 $\mathcal{O}$。}
\end{figure}

更详细地说，一个在 $\mathcal{T}$ 的 $G^{(p)}$ 对称性下带电荷
\begin{equation}
\widehat{g}\in\widehat{G}^{(p)}\,,
\tag{4.180}
\end{equation}
的算符 $\mathcal{O}$，处于 $\mathcal{T}/G^{(p)}$ 的
\begin{equation}
G^{(d-p-2)}=\widehat{G}^{(p)}
\tag{4.181}
\end{equation}
对称性的 $\widehat{g}$-扭曲扇区中。

\begin{examplebox}{例 4.14：从 $\SU(2)$ 到 $\SO(3)$：Wilson 线}
Wilson 线
\begin{equation}
\mathcal{W}_j\,,\qquad j\notin\Z\,,
\tag{4.182}
\end{equation}
在 $\SU(2)$ 理论的 $\Z_2$ 1-形式对称性下带电。正如我们先前讨论过的，$\SO(3)$ 理论中仅有的真 Wilson 线算符是 $j\in\Z$ 的 $\mathcal{W}_j$。上述所有 $j\notin\Z$ 的 Wilson 线都处于 $\SO(3)$ 理论 $\Z_2$ 1-形式对称性的扭曲扇区中，附着于磁 1-形式对称性的生成元 (3.113) 上。
\end{examplebox}

% [Tutor Note: p.92 例 4.14——$\SU(2)$ 的 Wilson 线 $\mathcal{W}_j$ 用自旋 $j$ 标记，$\SO(3)$ 只允许整数自旋，故 $j\notin\Z$ 的线在 $\SO(3)$ 中不再是真算符，而成为扭曲扇区算符，其端点必须接在 (3.113) 的磁 1-形式对称性生成元上。这与例 4.13 的 't Hooft 线构成「带电 $\leftrightarrow$ 扭曲」的配对，也是 (4.177) 与 (4.180) 的现实版本。]

\para{带电+扭曲变为扭曲+带电。}为简化论证，上面我们假设了扭曲扇区算符不带电、带电算符不在扭曲扇区中。我们可以放弃这一假设，考虑如图 19 所示既处于扭曲扇区又带电的算符。注意，只有当 $d$ 与 $p$ 满足
\begin{equation}
d=2p+2
\tag{4.183}
\end{equation}
的关系时这才可能。设 $\mathcal{O}$ 是 $\mathcal{T}$ 中具有下述性质的算符
\begin{equation}
(q,t)=(\widehat{g},g)\in\widehat{G}^{(p)}\times G^{(p)}\,,
\tag{4.184}
\end{equation}

% [Tutor Note: p.92 式 (4.183)——$d=2p+2$ 的来源可直接从维数读出：处于 $G^{(p)}$ 扭曲扇区的算符余维为 $p+2$，即维数 $d-p-2$；而对偶对称性 (4.181) 是 $G^{(d-p-2)}=\widehat{G}^{(p)}$。当 $d-p-2=p$，即 $d=2p+2$ 时，$p$-形式与 $(d-p-2)$-形式重合，同一个算符才能同时带「扭曲」和「带电」两种标签。核对锚点：$d=4,p=1$ 满足 $4=2\cdot 1+2$，正是本文反复使用的 4d $\SU(2)/\SO(3)$ 例子；$d=3,p=0$ 也满足 $3=2\cdot 0+2$，对应 3d 的 $\Z_2$ 0-形式对称性例子。]
